% Part: first-order-logic
% Chapter: sequent-calculus

\documentclass[../../../include/open-logic-chapter]{subfiles}

\begin{document}

\iftag{FOL}
      {\olchapter{fol}{seq}{حساب التتابعيات}}
      {\olchapter{pl}{seq}{حساب التتابعيات}}

\begin{editorial}
  موضوع هذا الفصل حساب غنتسن القياسي للتتابعيات، LK، في منطق الرتبة الأولى الكلاسيكي. وينبغي أن يُزاد فيه من الأمثلة والتمارين. وأما المادة التي تبحث حساب التتابعيات من حيث هو نظام برهان، فإدراجها واستبعادها بالوسم «prfLK».
\end{editorial}

\olimport{rules-and-proofs}

\olimport{propositional-rules}

\iftag{FOL}{%
  \olimport{quantifier-rules}
}{}

\olimport{structural-rules}

\olimport{derivations}

\olimport{proving-things}

\iftag{FOL}{%
  \olimport{proving-things-quant}
}{}

\olimport{proof-theoretic-notions}

\olimport{provability-consistency}

\olimport{provability-propositional}

\iftag{FOL}{%
  \olimport{provability-quantifiers}
}{}

\olimport{soundness}

\iftag{FOL}{%
  \olimport{identity}
  \olimport{soundness-identity}
}{}

\OLEndChapterHook

\end{document}
